\section{The Lifecycle as a Real Cycle}
\label{p11:sec:cycle}

\noindent
This section specifies the object the paper takes as primary. The lifecycle of a shared flourishing is a cycle, in the dynamical sense recalled in \xref{p11:sec:dynsys}: a sustained cyclic movement that maintains the relation by its continued turning, rather than a linear sequence of stages completed once and for all. The section establishes three things about this cycle: that it is not linear but a movement whose moments interpenetrate and feed back; that its persistence is the condition of the relation's existence, so that the cessation of the cycle is the relation's extinction; and that it turns, across the time of a relation, through nested repetitions from a beginning, through a marriage, into the later years.

\subsection{The moments of the cycle, and their non-linearity}
\label{p11:sec:cyc-moments}

The cycle is articulated into five moments: cognition, the apprehension of what the flourishing of this shared life consists in; practice, the cultivation and the action through which it is pursued; evaluation, the assessment of how the relation stands; reflection, the re-examination of that assessment and of the cognition that framed it; and re-cognition, the revised apprehension with which a further turn of the cycle begins. These are moments and not stages. A stage is completed and left behind; a moment is a recurring aspect of a movement that is always, in some measure, engaged in all of them at once. The parties do not first complete their cognition of the good, then practise, then evaluate, then reflect, then begin again. They practise on a cognition that is still forming; they evaluate continuously, as \xref{p11:sec:appraisal} established, while they practise; their reflection revises their practice mid-course; and their cognition is reconceived without waiting for a turn of the cycle to be complete. The moments interpenetrate and feed back upon one another, and the articulation into five is an analytic device for treating, in turn, aspects of a movement that does not in fact proceed in turn.

One moment in particular refuses localisation. Evaluation, as the basal and continuous appraisal of \xref{p11:sec:appraisal}, is not one station of the cycle but an undercurrent running through all of them: the parties are appraising the state of the relation while they cognise its good, while they practise, while they reflect, and while they reconceive. The treatment of evaluation as a distinct moment (\xref{p11:sec:evaluation}) concerns its explicit and instrumented form, the deliberate measurement undertaken at particular junctures; the basal evaluation, the unbroken belief updating, pervades the whole. The cycle is, in this respect, a movement with a continuous monitoring of its own state, in the manner of a system that maintains itself by never ceasing to estimate where it stands.

\subsection{The persistence of the cycle as the condition of the relation}
\label{p11:sec:cyc-persistence}

The central claim of the section concerns what the turning of the cycle is for. It is not that the cycle is a useful description of a relation that would exist in any case. The turning of the cycle is the activity by which the relation exists, and its cessation is the relation's end.

\begin{claim}[The persistence of the cycle is the condition of the relation's existence]
A relation is a living process maintained by the continued turning of its cycle, and not a state that, once attained, persists of itself. The continued generation of value, its circulation and return, the unbroken appraisal of the relation's state, and the practice taken upon that appraisal, are the activity by which the relation holds itself in being, in the manner of a self-maintaining system that exists only so long as it continues to infer and to act (\xref{p11:sec:ap-active}). A relation in which this activity ceases does not persist in a quiescent but intact condition; it decays toward the terminal rest in which value is no longer generated or returned and the parties no longer revise their sense of one another. The cessation of the cycle is therefore not a relation at rest but a relation extinguished. To be in a living relation is to be engaged in the continued turning of its cycle; to cease the turning is to let the relation lapse.
\end{claim}

\noindent
This claim distinguishes the present account from any treatment of relational well-being as a state to be reached and thereafter possessed. There is no state of a achieved flourishing that, once secured, maintains itself without the continued activity of the cycle, just as there is no fixed point of a living relation that is not its cessation (\xref{p11:sec:dyn-fixed}). The good of a shared life is not banked; it is generated, again and again, in the turning, and a relation that ceases to generate it does not hold what it had but loses it. The implication for practice is developed throughout the paper: the work of a shared life is not the attainment of a condition but the sustaining of a movement, and the most consequential failure is not the failure to reach a state but the cessation of the turning.

\subsection{The cycle across the time of a relation}
\label{p11:sec:cyc-time}

The cycle turns, and a relation is many turns of it, nested and accumulating. A single exchange may be a small turn, cognised, practised, appraised, reflected upon, and reconceived within an afternoon. A season of a relation is a larger turn, and the passage from courtship to a settled life, from a settled life through its trials, into the later years, is the largest turn the paper considers. These turns are nested: the small turns compose the seasons, and the seasons compose the arc of a shared life. The marriage at which a relation may arrive is not the end of the cycle but a transformation of its conditions, after which the same moments recur under altered circumstances, the cognition of the good revised by years of shared practice, the practice settled into the habits and household of \xref{p11:sec:practice}, the evaluation informed by a long history, the reflection deepened, and the re-cognition carrying forward what the prior turns have accumulated.

This accumulation across turns is the lifecycle's holonomy (\xref{p11:sec:dyn-phase}). The re-cognition with which a turn ends is not identical to the cognition with which it began; the difference, accumulated over a turn, is the phase the turn has acquired, and accumulated over many turns, it is the long ascent or the long decline of a shared life. A relation whose turns accumulate a positive phase rises across its history, its cognition of the good deepening, its practice enriching, its parties developing through their shared life into more than they were. A relation whose turns accumulate a negative phase declines across its history, however intact its form, depleting its parties turn by turn. And a relation whose cycle ceases to turn neither rises nor declines but ends. These three possibilities, the rising spiral, the declining cycle, and the cessation, are the three regimes the paper treats as a whole in \xref{p11:sec:regimes}, after it has traversed the moments of a single turn. The paper now enters the cycle at its first moment, the cognition of what the flourishing of a shared life consists in.
